% ID: SEL-D0081-OA-E04
% ESTADO: cotejada. Segunda pasada visual y borrador compilado revisados.
% Fuente DOC: Selectividad/sel_2012_matematicas_aplicadas/3 MATEMÁT. APLIC. CC. SS/EXAMEN 2-ANDALUCÍA.doc
% SHA-256 original: bcec57ae525fa5e23ce74631274a2abda0b5b16ac47ce6108ba8b5572d4aa634
% Página convertida: 1; opción A; ejercicio 4.

Se acepta que los rendimientos anuales, medidos en porcentajes, que producen los depósitos bancarios a plazo, se distribuyen según una ley Normal con desviación típica $1.8$ y se pretende realizar una estimación del rendimiento medio de los mismos. Para ello, se tiene una muestra de 36 entidades bancarias en las que se observa que el rendimiento medio de los depósitos es del $2.5$.
\begin{enumerate}[label=\alph*)]
\item \textbf{(1.5 puntos)} Calcule un intervalo de confianza, al $96\%$, para el rendimiento medio de los depósitos a plazo. \textquestiondown{}Cuál es el error máximo cometido en la estimación?
\item \textbf{(1 punto)} Manteniendo el mismo nivel de confianza, \textquestiondown{}cuál debe ser el tamaño mínimo de la muestra para estimar el rendimiento medio de los depósitos con un error máximo de $0.5$?
\end{enumerate}
